% Generated draft: requires mathematical and source review.
\documentclass[11pt]{article}
\usepackage[margin=25mm]{geometry}
\usepackage{fontspec}
\setmainfont{Latin Modern Roman}
\usepackage{amsmath,amssymb,mathtools,unicode-math}
\setmathfont{Latin Modern Math}
\usepackage{enumitem,needspace,xurl,hyperref}
\setlength{\parindent}{0pt}
\setlength{\parskip}{5pt}
\setlength{\emergencystretch}{3em}
\newcommand{\N}{\mathbb{N}}
\newcommand{\Z}{\mathbb{Z}}
\newcommand{\Q}{\mathbb{Q}}
\newcommand{\R}{\mathbb{R}}
\newcommand{\C}{\mathbb{C}}
\newcommand{\F}{\mathbb{F}}
\newcommand{\A}{\mathbb{A}}
\newcommand{\PP}{\mathbb{P}}
\newcommand{\E}{\mathbb{E}}
\newcommand{\eps}{\varepsilon}
\newcommand{\dd}{\,\mathrm{d}}
\begin{document}
\section*{UnsolvedMath ID: 30000233}
\textbf{Source problem code:} catalogue entry
\subsection{Definitions and mathematical statement}
Determine the final shape of the shrinking flat side in the Gauss-curvature flow of a weakly convex surface with flat sides.

In the usual interpretation for surfaces in $\R^3$, a smooth portion of an evolving surface $M_t$ has normal velocity
\[
 \partial_t X=-K\nu,
\]
where $\nu$ is the outward unit normal and $K=\kappa_1\kappa_2$ is Gauss curvature. Weak convexity means nonnegative principal curvatures where these are defined. A flat side is a two-dimensional planar region $F_t\subset M_t$ with nonempty relative interior. The equation on nonsmooth portions must be understood in the solution class intended by the report; that class is not specified here. The ambient space and compactness assumptions are also not explicitly supplied.

If the flat region collapses at a time $T$, its unrescaled limit need not encode its shape. A precise version of the question must specify a normalization and ask for a limit of normalized planar regions, for example regions of the form $a_t(F_t-b_t)$ after identifying the supporting plane with $\R^2$, with $b_t$ a prescribed center and $a_t>0$ a prescribed scale. The record does not prescribe these choices or the topology of convergence, so this notation describes the missing formulation rather than selecting it.

The dated literature assessment refers to a conjectured ellipse as the final shape. In an appropriate normalization, this would ask whether the limiting region is an ellipse (equivalently, whether its boundary is an ellipse), together with existence of the asserted limit. The supplied statement itself asks to determine the shape and does not specify the ellipse conjecture, the normalization, or whether the limit is unique. No particular area, diameter, affine, or other normalization is imposed in this draft, and no universal convergence theorem is asserted.
\subsection{Short English statement}
The question concerns the shape of a planar face as it disappears during Gauss-curvature flow. The supplied assessment mentions an elliptical limiting shape, but a precise conjecture still needs initial-data assumptions, a solution notion, a rescaling, and a meaning of convergence.
\subsection{Sources}
\begin{itemize}
\item Local UnsolvedMath registry, permanent ID 30000233, source code catalogue entry, ``Flat-Side Evolution Under Gauss Curvature Flow.'' Proposed year 2005; proposer unspecified.
\item \emph{Nonlinear Evolution Problems} (2006), Oberwolfach Reports, pp.\ 1371--1430, DOI 10.4171/owr/2005/25. Original Oberwolfach Report 25/2005: \url{https://doi.org/10.4171/owr/2005/25}. The displayed 2006 report year differs from the DOI and proposed year 2005.
\item Supplied reference: Caputo--Daskalopoulos, ``highly degenerate harmonic mean curvature flow,'' \url{https://arxiv.org/abs/0804.3936}. Its stated flow differs from the Gauss-curvature flow in the problem; its relevance requires source checking.
\item Supplied reference: Daskalopoulos, ``survey of degenerate curvature flows,'' \url{https://www.intlpress.com/site/pub/files/_fulltext/journals/sdg/2014/0019/0001/SDG-2014-0019-0001-a004.pdf}.
\item Statement status exact, with the supplied verification note ``clear and internally complete; retained byte-for-byte.'' The original question is reproduced above; its precision still depends on omitted report conventions.
\item Literature triage checked 2026-08-21: open. It reports known short-time regularity, interface motion, and collapse of the flat side, while saying later work does not give a general proof of the conjectured ellipse. These are supplied dated assessments, not independently verified results. Created 2026-08-25; updated 2026-09-16.
\item Registry curation: Ulam AI / UnsolvedMath Contributors, CC BY 4.0; underlying report and reference provenance remain as supplied. Unreviewed statement-only draft.
\end{itemize}
\end{document}
